
\begin{vidu}%[3P1-14.1K][Loai1] 
Một vật dao động điều hòa với biên độ 10 cm và chu kì 12 s.
\begin{itemize}
\item[1.] Quãng đường lớn nhất vật đi được trong $2\ s$ là \dotfill
\item[2.] Quãng đường lớn nhất vật đi được trong $3\ s$ là \dotfill
\item[3.] Quãng đường lớn nhất vật đi được trong $4\ s$ là \dotfill
\item[4.] Quãng đường nhỏ nhất vật đi được trong $2\ s$ là \dotfill
\item[5.] Quãng đường nhỏ nhất vật đi được trong $3\ s$ là \dotfill
\item[6.] Quãng đường nhỏ nhất vật đi được trong $4\ s$ là \dotfill
\item[7.] Quãng đường lớn nhất vật đi được trong $32\ s$ là\dotfill
\item[8.] Quãng đường lớn nhất vật đi được trong $33\ s$ là \dotfill
\item[9.] Quãng đường lớn nhất vật đi được trong $34\ s$ là \dotfill
\item[10.] Quãng đường nhỏ nhất vật đi được trong $38\ s$ là \dotfill
\item[11.] Quãng đường nhỏ nhất vật đi được trong $39\ s$ là \dotfill
\item[12.] Quãng đường nhỏ nhất vật đi được trong $40\ s$ là \dotfill
\item[13.] Thời gian ngắn nhất để vật đi được quãng đường $10\ cm$ là \dotfill
\item[14.] Thời gian dài nhất để vật đi được quãng đường $40-10\sqrt{3}\ cm$ là \dotfill
\end{itemize}
\loigiai{
\begin{itemize}
\item[1.] Quãng đường lớn nhất vật đi được trong $2\ s$ là 10 cm
\item[2.] Quãng đường lớn nhất vật đi được trong $3\ s$ là $10\sqrt{2}\ cm$
\item[3.] Quãng đường lớn nhất vật đi được trong $4\ s$ là $10\sqrt{3}\ cm$
\item[4.] Quãng đường nhỏ nhất vật đi được trong $2\ s$ là $20-10\sqrt{3}\ cm$
\item[5.] Quãng đường nhỏ nhất vật đi được trong $3\ s$ là $20-10\sqrt{2}\ cm$
\item[6.] Quãng đường nhỏ nhất vật đi được trong $4\ s$ là $10\ cm$
\item[7.] Quãng đường lớn nhất vật đi được trong $32\ s$ là $110\ cm$
\item[8.] Quãng đường lớn nhất vật đi được trong $33\ s$ là $100+10\sqrt{2}\ cm$
\item[9.] Quãng đường lớn nhất vật đi được trong $34\ s$ là $100+10\sqrt{3}\ cm$
\item[10.] Quãng đường nhỏ nhất vật đi được trong $38\ s$ là $140-10\sqrt{3}\ cm$
\item[11.] Quãng đường nhỏ nhất vật đi được trong $39\ s$ là $140-10\sqrt{2}\ cm$
\item[12.] Quãng đường nhỏ nhất vật đi được trong $40\ s$ là $130\ cm$
\item[13.] Khoảng thời gian ngắn nhất để vật đi được quãng đường $10\ cm$ là 2 s.
\item[14.] Khoảng thời gian dài nhất để vật đi được quãng đường $40-10\sqrt{3}\ cm$ là 8 s.
\end{itemize}
}
\end{vidu} 
%\dotlineEX{5}
%\loai{Quãng đường nhỏ nhất không tách thời gian}
\begin{vidu}%[3P1-14.1K][Loai1] 
Một chất điểm dao động điều hòa với biên độ 10 cm và chu kì 1,5 s. Trong khoảng thời gian 0,5 s, quãng đường nhỏ nhất mà vật đi được là
\choice
{$10\sqrt{3}$ cm}
{$10\sqrt{2}$ cm}
{\True 10 cm}
{20 cm}
\loigiai{
\begin{itemize}
\item Ta có: $\omega = \dfrac{2\pi}{T} = \dfrac{4\pi}{3}$ rad/s.
\item Trong khoảng thời gian $\Delta t = 0,5\ s$ chất điểm quay được một góc $\Delta \varphi = \dfrac{2\pi}{3}
\Rightarrow S_{\min} = 2.10(1-\cos\dfrac{\pi}{3}) = 10$ cm.
\end{itemize}
}
%\haidong{6}
\end{vidu} %\dotlineEX{5}
%\loai{Quãng đường lớn nhất cần tách thời gian}
\begin{vidu}%[3P1-14.1K][Loai1]
Một chất điểm dao động điều hòa với chu kì $T = 4$ s, biên độ $A = 4$ cm. Quãng đường lớn nhất mà vật đi được trong quãng thời gian 9 s là
\choice
{\True $32 + 4\sqrt{2}$ cm}
{$32 - 8\sqrt{2}$ cm}
{$40 - 4\sqrt{2}$ cm}
{$40 + 8\sqrt{2}$ cm}
\loigiai{
\begin{itemize}
\item Ta có: $\omega = \dfrac{2\pi}{T} = \dfrac{\pi}{2}$ rad/s.
\item Ta lại có: $t = 9\ s = 2.4 + 1\Rightarrow S = 2.4.4+ \Delta s = 32+ \Delta s$.
\item Trong khoảng thời gian $\Delta t = 1\ s$ chất điểm quay được một góc $\Delta \varphi = \dfrac{\pi}{2}
\Rightarrow \Delta s_{\max} = 2.4\sin\dfrac{\pi}{4} = 4\sqrt{2}\ cm \Rightarrow
S_{\max} = 32+4\sqrt{2}$ cm.
\end{itemize}
}
%\haidong{6}
\end{vidu} %\dotlineEX{4}
%\loai{Quãng đường nhỏ nhất cần tách thời gian}
\begin{vidu}%[3P1-14.1K][Loai1] 
Một chất điểm dao động điều hòa với biên độ $A = 6$ cm và chu kì $T = 1,5$ s. Quãng đường ngắn nhất vật đi được trong quãng thời gian 1 s là
\choice
{$24\sqrt{3}$ cm}
{\True $(24 - 6\sqrt{3})$ cm}
{$(24 + 6\sqrt{3})$ cm}
{$(12 + 6\sqrt{2})$ cm}
\loigiai{
\begin{itemize}
\item Ta có: $\omega = \dfrac{2\pi}{T} = \dfrac{4\pi}{3}$ rad/s.
\item Ta lại có: $t = 1\ s = 0,75 + 0,25\Rightarrow S = 2.6+ \Delta s = 12+ \Delta s$. 
\item Trong khoảng thời gian $\Delta t = 0,25\ s$ chất điểm quay được một góc $\Delta \varphi = \dfrac{\pi}{3}\Rightarrow
\Delta s_{\min} = 2.6(1-\cos\dfrac{\pi}{6}) = 12-6\sqrt{3}\ cm \Rightarrow
S_{\min} = 24-6\sqrt{3}$ cm.
\end{itemize}
}
%\haidong{8}
\end{vidu} %\dotlineEX{4}
%\loai{Quãng đường vật đi được có thể, không thể bằng}
\begin{vidu}%[3P1-14.1K][Loai1]
Một vật dao động điều hòa với phương trình $x=4\cos \left( \dfrac{4\pi}{3}t+\dfrac{\pi}{3} \right)$ cm trên trục Ox. Trong 1,75 s thì quãng đường đi được của vật \textbf{không} thể bằng
\choice
{18 cm}
{\True 17 cm}
{19 cm}
{20 cm}

\loigiai{
\begin{itemize}
\item $T = 1,5\ s \Rightarrow \Delta t  =1,75\ s=2.\dfrac{T}{2}+\dfrac{T}{6}\Rightarrow S_{\max}=4A+2A\sin \dfrac{\pi}{6}=20$ cm;  $S_{\min}=4A+2A\left( 1-\cos \dfrac{\pi}{6} \right)=17,07$ cm.
\item Quãng đường vật S có thể đi được phải thỏa mãn $S_{\min} < S < S_{\max}$.
\item  Trong 1,75 s thì quãng đường đi được của vật không thể bằng 17 cm.
\end{itemize}
}
%\haidong{11}
\end{vidu} %\dotlineEX{4}
%\loai{Tỉ số $S_{\max}$, $S_{\min}$}
\begin{vidu}%[3P1-14.1K][Loai1]
Một vật dao động điều hoà dọc theo trục Ox, quanh vị trí cân bằng O với biên độ A và chu kì T. Trong khoảng thời gian $\dfrac{T}{6}$, tỉ số quãng đường lớn nhất, nhỏ nhất mà vật có thể đi được là
\choice
{2}
{\True $2+\sqrt{3}$}
{$2+\sqrt{2}$}
{3}
\loigiai{
\begin{itemize}
\item Tỉ số: $\Delta t = \dfrac{T}{6} \Rightarrow \dfrac {S_{\max}}{S_{\min}} = \dfrac{2A\sin \dfrac{\pi \Delta t}{T}}
{2A\left( 1 - \cos \dfrac{\pi \Delta t}{T} \right)} = \dfrac{2A\sin \dfrac{\pi}{6}}{2A\left( {1 - \cos \dfrac{\pi}{6}} \right)} = 2 + \sqrt 3 $.
\item 
\end{itemize}}
%\haidong{8}
\end{vidu} %\dotlineEX{4}

%\loai{Thời gian ngắn nhất để vật đi quãng đường $S<2A$}
\begin{vidu}%[3P1-14.2K][Loai1]
Một vật dao động điều hòa với biên độ bằng 4 cm, chu kì 2 s. Khoảng thời gian nhỏ nhất vật cần để đi được quãng đường  $4\sqrt{3}$  cm là
\choice
{$\dfrac{1}{3}\ s$}
{$\dfrac{1}{6}\ s$}
{\True $\dfrac{2}{3}\ s$}
{$\dfrac{3}{4}\ s$}
\loigiai{
\begin{itemize}
\item Ta có: $S = 4\sqrt 3 < 2A = 8 \Rightarrow 4\sqrt 3 = 2A\sin \dfrac{\pi \Delta t_{\min}}{T} = 2.4.\sin \dfrac{\pi \Delta t_{\min }}{T} \Rightarrow \Delta t_{\min} = \dfrac{T}{3} = \dfrac{2}{3}\ s$.

\end{itemize}
}
%\haidong{8}
\end{vidu} %\dotlineEX{6}
%\loai{Thời gian dài nhất để vật đi quãng đường $S<2A$}
\begin{vidu}%[3P1-14.2K][Loai1]
Một vật dao động điều hoà với tần số 0,5 Hz, biên độ A. Khoảng thời gian lớn nhất để vật đi được quãng đường bằng A là
\choice
{1 s}
{$\dfrac{1}{3}$ s}
{$\dfrac{1}{2}$ s}
{\True $\dfrac{2}{3}$ s}
\loigiai{
\begin{itemize}
\item Ta có: $f = 0,5\ Hz \Rightarrow T = 2$ s.
\item Khoảng thời gian để vật đi được quãng đường bằng A lớn nhất khi góc quay $\Delta \varphi$ lớn nhất\\
$\Rightarrow \Delta \varphi_{\max} = \dfrac{2\pi}{3}\Rightarrow t = \dfrac{T}{3} = \dfrac{2}{3}\ s$.

\end{itemize}
}
%\haidong{8}
\end{vidu} %\dotlineEX{7}
%\loai{Thời gian ngắn nhất để vật đi quãng đường $S>2A$}
\begin{vidu}%[3P1-14.2K][Loai1]
Một vật dao động điều hòa với biên độ bằng 4 cm. Khoảng thời gian nhỏ nhất vật cần để đi được quãng đường 12 cm là 0,8 s. Số dao động toàn phần mà vật thực hiện trong khoảng thời gian mỗi phút là
\choice
{45}
{43}
{34}
{\True 50}

\loigiai{
\begin{itemize}
\item Ta có: $S=12=2A+A$
\item Để đi 2A luôn cần $\dfrac{T}{2}$ , thời gian nhỏ nhất để đi A là:
\begin{align*}
A=2A\sin \dfrac{\pi \Delta t_{\min}}{T}\Rightarrow \Delta t_{\min}=\dfrac{T}{6}.
\end{align*}
\item Vậy thời gian ngắn nhất để đi 12 cm là:
$\Delta t=\dfrac{T}{2}+\dfrac{T}{6}=\dfrac{2T}{3}=0,8\ s\Rightarrow T = 1,2\ s.$
\item Trong một phút số dao động toàn phần vật thực hiện là: $\dfrac{60}{T}=50$.
\end{itemize}
}
%\haidong{8}
\end{vidu} %\dotlineEX{7}
